% Generated draft: requires mathematical and source review.
\documentclass[11pt]{article}
\usepackage[margin=25mm]{geometry}
\usepackage{fontspec}
\setmainfont{Latin Modern Roman}
\usepackage{amsmath,amssymb,mathtools,unicode-math}
\setmathfont{Latin Modern Math}
\usepackage{enumitem,needspace,xurl,hyperref}
\setlength{\parindent}{0pt}
\setlength{\parskip}{5pt}
\setlength{\emergencystretch}{3em}
\newcommand{\N}{\mathbb{N}}
\newcommand{\Z}{\mathbb{Z}}
\newcommand{\Q}{\mathbb{Q}}
\newcommand{\R}{\mathbb{R}}
\newcommand{\C}{\mathbb{C}}
\newcommand{\F}{\mathbb{F}}
\newcommand{\A}{\mathbb{A}}
\newcommand{\PP}{\mathbb{P}}
\newcommand{\E}{\mathbb{E}}
\newcommand{\eps}{\varepsilon}
\newcommand{\dd}{\,\mathrm{d}}
\begin{document}
\section*{UnsolvedMath ID: 30003631}
\textbf{Source problem code:} catalogue entry
\subsection{Definitions and mathematical statement}
The supplied question does not define connected monodromy. The following makes explicit the usual interpretation in terms of algebraic monodromy on rational first cohomology; this interpretation requires confirmation against the original report.

A Kodaira fibration is understood here to be a smooth surjective holomorphic map \(f\colon S\to B\) with connected fibres, where \(S\) is a smooth connected projective complex surface, \(B\) is a smooth connected projective complex curve of genus \(b\geq 2\), and the fibres are smooth curves of genus \(g\geq 2\) whose complex structures are not all isomorphic. For a base point \(t\in B\), put \(F_t=f^{-1}(t)\) and \(V=H^1(F_t,\Q)\). The cup product gives a nondegenerate alternating form \(\langle\ ,\ \rangle\) on \(V\). Parallel transport defines a monodromy representation
\[
\rho_f\colon\pi_1(B,t)\longrightarrow\operatorname{Sp}(V,\langle\ ,\ \rangle).
\]
Under the stated interpretation, define the connected monodromy group by
\[
G_f=\left(\overline{\rho_f(\pi_1(B,t))}^{\,\mathrm{Zar},\Q}\right)^\circ,
\]
where the closure is the algebraic Zariski closure over \(\Q\), and the superscript \(\circ\) denotes its identity component.

\textbf{Classification question.} Which connected algebraic groups occur as \(G_f\) for Kodaira fibrations \(f\colon S\to B\), allowing the base and fibre genera to vary? More explicitly, for each \(g\geq 2\), determine which connected algebraic \(\Q\)-subgroups of \(\operatorname{Sp}_{2g,\Q}\) arise in this way, up to conjugacy by a rational symplectic change of basis, as \(b\geq 2\) and \(f\) vary.

The subgroup formulation specifies an ambient representation and an equivalence relation beyond those stated in the source record. The record does not settle whether the intended classification is instead up to abstract algebraic-group isomorphism, after extension of scalars, or under another convention for connected monodromy. It supplies no restriction to a particular fibre genus. Its literature assessment refers to a specified genus-three construction without supplying that construction or its hypotheses; those qualifications cannot be incorporated as a precise subsidiary statement.
\subsection{Short English statement}
Classify the connected monodromy groups arising from Kodaira fibrations. Under the explicitly stated cohomological interpretation, these are identity components of rational Zariski closures of monodromy on the first cohomology of a fibre. The intended field and classification equivalence require confirmation. The supplied record labels the problem partially solved, but this draft does not independently verify that assessment.
\subsection{Sources}
\begin{itemize}
\item Local registry: \texttt{data/unsolvedmath/problems.json}, permanent UnsolvedMath ID \(30003631\), source problem code \texttt{catalogue entry}, title \emph{Connected Monodromy Groups of Kodaira Fibrations}. The supplied record gives proposed year 2017 and identifies no proposer.
\item Supplied report citation: \emph{Mini-Workshop: Interactions between Low-dimensional Topology and Complex Algebraic Geometry (2018)}, Oberwolfach Reports, pp. 2987--3033, DOI \texttt{10.4171/owr/2017/50}. \url{https://doi.org/10.4171/owr/2017/50}. The supplied literature-source label is \emph{Oberwolfach Report 50/2017}; the relationship between the citation's 2018 date and the report's 2017 designation has not been verified.
\item Supplied literature source: Flapan, \url{https://doi.org/10.1002/mana.202000189}. No article title, full author name, or theorem reference is supplied. The supplied assessment says that Flapan classifies possible connected monodromy groups in a specified genus-three construction, while realizability in remaining cases and the question for arbitrary fibre genus are not completely settled. The construction and the remaining cases are not specified in the record.
\item Literature assessment date: \texttt{2026-08-21}. The status \texttt{partially\_solved} and the preceding assessment are supplied dated curation derived from the OWR record, not independently verified literature findings. The record explicitly requests independent verification of time-sensitive status and cited claims.
\item The supplied statement status is \texttt{corrected\_verified}, with the explanation that mathematical content was preserved while notation, typography, and sentence structure were repaired for publication. This is source metadata, not an independent mathematical verification of this draft.
\end{itemize}
\end{document}
